\documentclass[preprint,11pt]{elsarticle}

\usepackage{graphicx}
\usepackage{booktabs}
\usepackage{amsmath}
\usepackage{siunitx}
\usepackage{xcolor}
\usepackage[hidelinks]{hyperref}
\usepackage{url}

\journal{Biomedical Signal Processing and Control}

\begin{document}

\begin{frontmatter}

\title{Which representation clusters? A chance-corrected benchmark of ECG
representations for unsupervised arrhythmia structure}

\author[uh]{Dongho Chun}
\author[af]{Axel Faes\corref{cor1}}
\cortext[cor1]{Corresponding author.}
\ead{axel.faes@uhasselt.be}
\affiliation[uh]{organization={Data Science Institute (DSI), Hasselt University}, city={Hasselt}, country={Belgium}}
\affiliation[af]{organization={Biomedical Data Sciences, Data Science Institute (DSI) and Biomedical Research Institute (BIOMED), Hasselt University}, city={Hasselt}, country={Belgium}}

\begin{abstract}
Unsupervised analysis of the electrocardiogram (ECG) promises label-free
discovery of cardiac structure, yet whether clustering succeeds depends less on
the algorithm than on how the signal is represented. We benchmark four
representation families for unsupervised separation of sinus rhythm (SR) from
atrial fibrillation (AF, including flutter) on PTB-XL: whole-recording raw flattening, a convolutional
reconstruction autoencoder, compact hand-crafted rhythm and morphology features,
and two frozen foundation-model embeddings (HuBERT-ECG and ECGFounder).
Evaluation is chance-corrected (adjusted Rand index, ARI), patient-grouped,
repeated over ten seeds, and bounded by a supervised ceiling. ECGFounder clusters
best (ARI up to 0.885), but 16 interpretable hand-crafted features are competitive
(0.711) and far ahead of the raw and autoencoder pipelines (ARI $\approx 0$), which
match a majority-class assignment despite supervised ceilings above chance.
The same ordering holds on a MIT-BIH beat-morphology task. Two cautions emerge: the
silhouette score rewards the degenerate one-cluster solution (silhouette 0.81 at
ARI 0) and cannot validate ECG clustering alone, and feature conditioning materially
changes the outcome. The ranking is largely expected, since SR and AF differ by
definition in rhythm regularity. The contribution is an evaluation protocol whose
chance-corrected controls locate where naive pipelines fail. On one 9:1
class-imbalanced subsample (one $k$-means run) the hand-crafted features overtake
ECGFounder (ARI 0.232 versus 0.060).
Under simulated federation across ten sites with unequal class proportions, ECGFounder
keeps the higher mean ARI (0.612 versus 0.591) but retains less of its centralised
ARI. On these two binary tasks, representation mattered more than clustering algorithm.
\end{abstract}

\begin{keyword}
electrocardiography \sep unsupervised clustering \sep representation learning \sep
foundation models \sep adjusted Rand index \sep PTB-XL
\end{keyword}

\end{frontmatter}

\section{Introduction}\label{sec:intro}

Most machine-learning work on the electrocardiogram (ECG) is supervised, but the
clinically interesting questions are often unsupervised: which patients group
together, what subtypes exist, and whether structure can be found without labels
that are expensive or unavailable. A natural first attempt is to cluster recordings
directly. Whether that attempt succeeds is usually attributed to the clustering
algorithm, yet the choice that dominates is the representation fed to it.

This study began as a re-examination of a master's project that clustered whole
12-lead recordings. That pipeline flattened each strip into a long vector, compressed
it with a dense autoencoder, and clustered the code. It reported high cluster purity
and a silhouette score near 0.89, but the purity simply equalled the majority-class
proportion in every setting, and the silhouette was measured on a degenerate solution
in which almost all points fell into a single cluster. The honest conclusion of that
work, that the difficulty lies in processing the raw signal rather than in the
clustering step, is correct and is the starting point here. We make it precise.

A representation can carry information about a label and still fail to cluster by it.
Linear separability and clusterability are distinct properties: a representation may
place the discriminative direction far from its dominant axes of variation, so that a
supervised probe recovers the label while an unsupervised partition does not. This
distinction has been reported in the self-supervised learning literature, where network
layers with nearly the same linear-probe accuracy differ in clustering quality
\cite{sslclustering2025}, and we treat it as background rather than as a finding. What
has not been measured, for ECG, is how the common representation families compare on
this axis under proper chance-corrected metrics.

The most directly related work benchmarks several ECG foundation models against one
another, including a cluster-separability analysis of embeddings pooled from four
12-lead datasets, one of which is PTB-XL \cite{lookingbeyond2026}.
It does not include a comparison against
hand-crafted features, and it omits the Massachusetts Institute of Technology and Beth
Israel Hospital (MIT-BIH) Arrhythmia Database. Those two gaps are exactly
where this benchmark adds value, because hand-crafted features remain the
interpretable and inexpensive default in clinical signal processing, and MIT-BIH is the
canonical morphology benchmark.

We therefore benchmark four representation families for unsupervised separation of
sinus rhythm (SR) from atrial fibrillation (AF), with the controls the original pipeline
lacked. The resulting ranking is not the contribution. Because SR and AF differ by
definition in the regularity of the rhythm, that compact hand-crafted features and a
large foundation model both recover the split is expected. What the benchmark adds is
what the controls reveal and how the representations behave under deployment stress.
The contribution is an evaluation protocol for unsupervised ECG representations, applied
to five representations on PTB-XL and three on MIT-BIH. (i) The protocol places raw, reconstruction-autoencoder,
hand-crafted, and two frozen foundation-model representations in one chance-corrected,
patient-grouped, ten-seed, significance-tested clustering benchmark, and bounds every
clustering score by a supervised ceiling computed under the same patient grouping, so that
a low score can be attributed either to geometry or to missing information. The ordering
is checked across the assumed number of clusters. (ii) Cluster-size checks and a sweep
over feature conditioners yield two reproducible cautions that explain the original
project's misleading numbers: internal validity indices reward the degenerate collapse,
and feature conditioning changes the result by a wide margin. (iii) Clusterability is
tested under two simulated deployment stressors. On one 9:1 class-imbalanced subsample
(one $k$-means run) the 16 hand-crafted features overtake ECGFounder (adjusted Rand index, ARI, 0.232 against
0.060). Under simulated
federation with unequal class proportions across sites, ECGFounder keeps the higher mean
ARI (0.612 against 0.591), while the hand-crafted features keep a larger share of their
centralised ARI (0.918 against 0.692) and vary less across site partitions. (iv) A
generalisation check on a MIT-BIH beat-morphology task, the dataset the closest prior
benchmark omits. We do not propose a new representation or clustering algorithm, we scope
every foundation-model statement to the two models we test, and we credit the
separability-versus-clusterability phenomenon to prior work.

\section{Related work}\label{sec:related}

\paragraph{Unsupervised ECG and representations.} Hand-crafted features built from the
irregularity of the RR interval (the interval between successive R peaks) and from beat
morphology encode the defining signatures of common
arrhythmias and underpin decades of automatic analysis \cite{hrvstandards1996,dechazal2004}.
Tensor and wavelet decompositions have been combined with spectral clustering for
unsupervised 12-lead classification \cite{wavelettensorecg2018}. Reconstruction
autoencoders optimise signal energy rather than class structure, which makes them
strong compressors but weak organisers of class geometry. Contrastive self-supervised
methods built on physiologically motivated augmentations and patient- or lead-level
invariances, including CLOCS \cite{clocs2021}, 3KG \cite{threekg2021}, and contrastive
predictive coding for 12-lead ECG \cite{mehari2022ssl}, established the pretraining
recipes that later foundation models scaled up. Self-supervised and foundation models for
ECG, such as HuBERT-ECG \cite{hubertecg2024} and ECGFounder \cite{ecgfounder2024}, learn
general representations from millions of recordings and are now evaluated as frozen
feature extractors. The wavelet-tensor and spectral-clustering pipeline of
\cite{wavelettensorecg2018} addresses a related unsupervised goal with a different
(tensor) representation. We focus on the four families above, which span the energy-only
to definition-encoding spectrum, and leave a tensor arm to future work.

\paragraph{Clustering evaluation.} Cluster purity is not chance-corrected: assigning
every point to the majority class attains a purity equal to that class's prevalence.
External, chance-corrected indices, the adjusted Rand index (ARI) \cite{hubertarabie1985} and
adjusted mutual information (AMI) \cite{vinh2010}, are near zero for such trivial solutions
and are the appropriate headline metrics. Internal indices such as the silhouette
\cite{rousseeuw1987} measure compactness in the embedding being clustered and can be
maximised by degenerate solutions, so they require cluster-size sanity checks.

\paragraph{Closest benchmark.} The ECG foundation-model benchmark of
\cite{lookingbeyond2026} reports ARI-based cluster separability of foundation-model
embeddings pooled from four 12-lead datasets, one of which is PTB-XL, and finds that
some models cluster by clinical label while others cluster by dataset. It compares foundation models to one another and states that prior
work omits a representation-wise evaluation. It does not add a hand-crafted baseline,
and it does not include MIT-BIH. This study adds both, and pairs each clustering score
with a supervised ceiling computed on the same features and patient grouping, so that a
low clustering score can be attributed to geometry rather than to missing information.
The protocol applies the supervised ceiling, cluster-size checks, a conditioner sweep
and two simulated deployment stressors to the non-foundation and the foundation
representations alike, so that both kinds of representation are judged under the same
controls on the same patient-grouped cohort.

\section{Methods}\label{sec:methods}

\subsection{Data and cohorts}

\paragraph{PTB-XL.} We use PTB-XL \cite{ptbxl2020,ptbxl2022physionet}, a public 12-lead database
(21\,799 records from 18\,869 patients, 10\,s strips). Rhythm labels are taken from the
rhythm statements of the Standard Communication Protocol for Computer-Assisted
Electrocardiography (SCP-ECG) that accompany each record. The SR-versus-AF cohort uses as
positives all records carrying the statement AFIB (atrial fibrillation) or AFLT (atrial
flutter), and as controls the records carrying SR without either statement. The positive
class therefore contains atrial flutter as well as atrial fibrillation
(\textcolor{red}{[AUTHOR: number of AFLT records among the 1\,260 positives]}), and we call
it AF throughout. PTB-XL encodes rhythm statements at a likelihood of zero meaning
``present'', so a presence-based rule is required. A confidence threshold would discard
every SR record. We draw a balanced, patient-disjoint cohort of 2\,542 records (1\,282 SR,
1\,260 AF), so the majority-class baseline is 0.504. The target size per class is the
number of distinct AF patients. SR records are shuffled with a fixed seed and taken one
per patient, and AF records are then shuffled and taken one per patient, skipping
patients already drawn. Lead use differs by representation (Section~\ref{sec:repr}): the
hand-crafted features are computed from lead II alone, whereas raw flattening, the
convolutional autoencoder and both foundation models receive all twelve leads. One lead
suffices for the RR-interval features, because R-peak timing is common to all leads, and
AF detection from a single lead is an established task \cite{clifford2017af}. The P-wave
and QRS-complex features, however, are then read from lead II only, so lead-specific morphology in
the other eleven leads is not used, and the effect of choosing another lead was not
evaluated. The hand-crafted representation therefore receives less input than the other
PTB-XL representations, and the comparison does not favour it in input. All evaluation is
patient-grouped: no patient appears in more than one fold or in both classes.

\paragraph{Signal input and preprocessing.} PTB-XL records (version 1.0.3) are read with
the \texttt{wfdb} Python package into 32-bit floating-point arrays of the physical signal
values in the units declared in each record header, without unit conversion. The main
benchmark uses the 100\,Hz versions of the records (1\,000 samples per lead). MIT-BIH
records are read the same way at their native sampling rate of 360\,Hz. No band-pass
filter or baseline-wander removal is applied to the inputs of raw flattening, the
convolutional autoencoder, the foundation models or the MIT-BIH representations. The
hand-crafted PTB-XL features are computed on lead II after the default cleaning step of
NeuroKit2. Resampling and normalisation per representation are listed in
Table~\ref{tab:repr}. After extraction, every representation passes through the feature
conditioner of Section~\ref{sec:clust} before clustering.

\paragraph{MIT-BIH.} For generalisation we use the MIT-BIH Arrhythmia Database
\cite{mitbih2001,physionet2000}. Each annotated beat is segmented into a window of
$[-300, +250]$\,ms around the R-peak (198 samples at 360\,Hz) across both recorded
channels. Beats whose window crosses the record boundary, and the first and last annotated
beat of each record, are skipped. Beats annotated N (normal) form one class and beats
annotated V (premature ventricular contraction) the other. At most 400 beats per record
and class are kept, from which 1\,500 per class are drawn with a fixed seed. This gives a
balanced normal-versus-premature-ventricular (N-vs-V) task of 3\,000 beats from 47
records. The record identifier is the grouping unit for cross-validation, and records are
not mapped to subjects. This is a morphology task and a deliberate contrast to the PTB-XL
rhythm task.

\subsection{Representations}\label{sec:repr}

We compare four families on PTB-XL (Table~\ref{tab:main}), and Table~\ref{tab:repr} gives
the input of each representation. All inputs are 32-bit floating-point arrays, and every
representation outputs a real-valued vector of the dimensionality given in the table.
The inputs differ by design: raw flattening and the convolutional autoencoder reproduce
the multi-lead input of the original pipeline, the foundation models receive 12-lead
input of the length, sampling rate and layout used in their pretraining, without the
filtering and, for HuBERT-ECG, the amplitude scaling that their authors apply (see below),
and the hand-crafted features use one rhythm lead. The comparison is therefore between representations as they would be deployed,
not under equal input.
\begin{itemize}
\item \textbf{Raw flattening} reproduces the original pipeline: each of the 12 leads is
downsampled by a non-overlapping moving average of window 6 (from 1\,000 to 166 samples)
and the leads are concatenated into one vector of 1\,992 dimensions, with no filtering or
normalisation. It carries amplitude and energy but no cross-recording temporal alignment.
\item \textbf{Convolutional autoencoder} (conv-AE): a one-dimensional (temporal)
convolutional autoencoder with a 32-unit bottleneck, trained to minimise the mean squared
reconstruction error. Its code keeps what is needed to reconstruct the signal, which need
not be the class structure. Each record is resampled with the fast Fourier transform (FFT)
from 1\,000 to 2\,500 samples per lead and z-scored per lead over time. Training stops
early on a 10\% validation split (architecture and training settings in
Table~\ref{tab:impl}). This is a properly specified reconstruction model, in contrast to
the original collapsed dense autoencoder.
\item \textbf{Hand-crafted} (16 features), computed from lead II with NeuroKit2
\cite{neurokit2021}. They encode the clinical definitions of rhythm irregularity and beat
morphology \cite{hrvstandards1996,dechazal2004}. Eight come from the detected R peaks:
the mean RR interval, the standard deviation of the RR intervals (SDNN), the root mean
square of successive RR differences (RMSSD), the fraction of successive RR differences
larger than 50\,ms (pNN50), the coefficient of variation of the RR intervals, the Shannon
entropy of a 16-bin RR histogram over 200 to 2\,000\,ms, the mean heart rate, and the
number of detected beats. Eight describe morphology after wavelet-based delineation: the
fraction of beats with a detected P wave, the mean and standard deviation of the P-wave
amplitude, the mean width of the QRS complex, the mean and standard deviation of the
R-peak amplitude, and the mean and standard deviation of the correlation of each beat
($\pm 250$\,ms) with the median beat of the record. All are continuous except the beat
count. A feature that cannot be computed for a record (for example, when no P wave is
detected) is replaced by the median of that feature over the cohort.
\item \textbf{Foundation embeddings}, frozen and used without fine-tuning. HuBERT-ECG
\cite{hubertecg2024} is a self-supervised transformer. It receives the first 5\,s of each
100\,Hz record (500 samples per lead), the 12 leads flattened lead after lead into one
6\,000-sample vector without filtering or normalisation, and the embedding is the
final-layer hidden state averaged over time (768 dimensions). The preprocessing code
released by its authors (Code and data availability) band-pass filters each record from
0.05 to 47\,Hz, resamples it to 500\,Hz and scales each lead to $[-1,1]$. For pretraining
it takes the first 5\,s, flattens the 12 leads lead after lead and decimates the sequence
by a factor of 5, which gives the same 6\,000-sample layout at 100\,Hz as used here. This
study applies neither the band-pass filter nor the scaling to $[-1,1]$, and it takes the
100\,Hz records distributed with PTB-XL instead of decimating the 500\,Hz records.
ECGFounder \cite{ecgfounder2024} is a convolutional network trained on over 10 million
ECGs annotated by cardiology experts with 150 label categories. It receives the full
record FFT-resampled to 5\,000 samples per lead and z-scored once over the whole 12-lead
array, and the embedding is the time average of the last convolutional stage, taken before
the classification layer (1\,024 dimensions). In the preprint version of
\cite{ecgfounder2024} its authors resample each recording to 500\,Hz by linear
interpolation, apply a 1\,Hz high-pass filter, a second-order 30\,Hz Butterworth low-pass
filter and a 50/60\,Hz notch filter, cut 10\,s segments and z-score each segment with its
own mean and standard deviation. This study uses the same $12 \times 5\,000$ input and the
same z-score, but none of the three filters, and FFT resampling of the 100\,Hz records in
place of linear interpolation (Section~\ref{sec:limitations}). The model definition in the
authors' code repository builds the network with batch normalisation and dropout disabled,
as this study does.
\end{itemize}
On MIT-BIH the rhythm features do not transfer to single beats, so we compare a raw
beat window (both channels, 396 dimensions), 21 beat-morphology and RR-context features,
and the conv-AE. The 21 features are five RR-context features (the previous and the next
RR interval, their ratio, the ratio of the previous to the local RR interval, and the
local RR interval, defined as the mean of up to six preceding intervals) and, for each of
the two channels, the signed peak amplitude (the sample of largest magnitude), the signal
energy, a QRS width count (the number of samples whose magnitude exceeds half the peak
magnitude), the maximum, the minimum, the
skewness, the kurtosis, and the correlation with the median beat of the 3\,000 analysed
beats. All are continuous except the width count. The foundation models were not applied
to MIT-BIH, because their input in this study is a 5\,s (HuBERT-ECG) or 10\,s (ECGFounder)
12-lead strip, whereas a beat window spans 550\,ms of two channels.

\begin{table}[t]
\centering
\footnotesize
\caption{Input of each representation. Leads: number of leads or channels used, or the
name of the single lead. Input: array per record or beat after resampling (FFT: resampled
with the fast Fourier transform). Filter / norm.: filtering or normalisation applied
before the representation is computed. All arrays are 32-bit floating point. Output:
dimensionality of the representation before conditioning. $^{a}$First 5\,s of the
100\,Hz record, 12 leads flattened lead after lead. $^{b}$Default cleaning filter of
NeuroKit2.}
\label{tab:repr}
\setlength{\tabcolsep}{4pt}
\begin{tabular}{lllllr}
\toprule
Representation & Leads & Input & Filter / norm. & Training & Output \\
\midrule
\multicolumn{6}{l}{PTB-XL, 100\,Hz records} \\
raw-flatten     & 12 & $12 \times 1000$ & none & none & 1\,992 \\
conv-AE         & 12 & $12 \times 2500$ (FFT) & z-score/lead & trained & 32 \\
hand-crafted    & II & $1000$ & NeuroKit2$^{b}$ & none & 16 \\
HuBERT-ECG      & 12 & $6000^{a}$ & none & frozen & 768 \\
ECGFounder      & 12 & $12 \times 5000$ (FFT) & z-score/record & frozen & 1\,024 \\
\midrule
\multicolumn{6}{l}{MIT-BIH, 360\,Hz records} \\
raw-beat        & 2  & $2 \times 198$ & none & none & 396 \\
beat-morphology & 2  & $2 \times 198$ & none & none & 21 \\
conv-AE         & 2  & $2 \times 2500$ (FFT) & z-score/channel & trained & 32 \\
\bottomrule
\end{tabular}
\end{table}

\subsection{Clustering and evaluation}\label{sec:clust}

Each representation is clustered with $k$-means, Ward agglomerative clustering, spectral
clustering, a Gaussian mixture model (GMM), and hierarchical density-based spatial
clustering of applications with noise (HDBSCAN) \cite{hdbscan2013}, as implemented in
scikit-learn \cite{sklearn2011}. For the first four the number of clusters is set to the
number of classes, $K=2$. HDBSCAN chooses its own number of clusters and labels points
that belong to no dense region as noise. Its minimum cluster size is
$\max(10, \lfloor n/50 \rfloor)$, which is 50 for the 2\,542 PTB-XL records and 60 for
the 3\,000 MIT-BIH beats, and its other parameters, including the minimum number of
samples, are library defaults. The noise label is passed unchanged to the ARI, the AMI,
the normalised mutual information and the Fowlkes-Mallows index, so all noise points
together form one additional cluster in these indices. Noise points are neither
discarded nor reassigned. This penalises HDBSCAN when most points are noise. The settings
of all clusterers are listed in Table~\ref{tab:impl}. No setting was tuned, and every
setting not listed is the scikit-learn default.

\paragraph{Metrics.} Let $n$ items be partitioned into clusters $U_1,\dots,U_R$ and
classes $V_1,\dots,V_C$, with contingency counts $n_{ij}=|U_i \cap V_j|$, cluster sizes
$a_i$ and class sizes $b_j$. The adjusted Rand index \cite{hubertarabie1985} counts the
pairs of items on which the clustering and the labels agree, corrected for chance:
\begin{equation}
\mathrm{ARI}=\frac{\sum_{ij}\binom{n_{ij}}{2}-E}
{\frac{1}{2}\left[\sum_i\binom{a_i}{2}+\sum_j\binom{b_j}{2}\right]-E},
\qquad
E=\frac{\sum_i\binom{a_i}{2}\sum_j\binom{b_j}{2}}{\binom{n}{2}}.
\end{equation}
Its expected value under random labelling is 0, its maximum is 1 for identical partitions,
and it can be negative. Adjusted mutual information \cite{vinh2010} measures the
information the clustering shares with the labels, corrected for chance,
$\mathrm{AMI}=(I-E[I])/(\frac{1}{2}[H(U)+H(V)]-E[I])$, where $I$ is the mutual
information of the two partitions, $H$ the entropy and $E[I]$ the expected mutual
information under random labelling. Its expected value is 0 and its maximum 1. The
normalised mutual information (NMI) is $I/(\frac{1}{2}[H(U)+H(V)])$, in $[0,1]$ and not
corrected for chance \cite{vinh2010}. The Fowlkes-Mallows index \cite{fowlkes1983} is the
geometric mean of pair precision and pair recall, in $[0,1]$. Purity,
$\frac{1}{n}\sum_i \max_j n_{ij}$ with noise points counted as misassigned, lies in
$[0,1]$ and equals the majority-class proportion for a single cluster. The silhouette
\cite{rousseeuw1987} of point $i$ is
$s(i)=(\bar d_{\mathrm{out}}(i)-\bar d_{\mathrm{in}}(i))/\max\{\bar d_{\mathrm{in}}(i),\bar d_{\mathrm{out}}(i)\}$,
where $\bar d_{\mathrm{in}}(i)$ is the mean distance from $i$ to the other points of its
cluster and $\bar d_{\mathrm{out}}(i)$ the mean distance to the points of the nearest other
cluster. It is averaged over points, lies in $[-1,1]$, is computed in the conditioned
representation without labels, and excludes HDBSCAN noise points. The largest-cluster
fraction is the size of the largest non-noise cluster divided by $n$. Without noise it
lies in $[1/K,1]$, and 1 marks a one-cluster solution. ARI and AMI are the headline
metrics. NMI, the Fowlkes-Mallows index, purity and the largest-cluster fraction of every
cell are given in Table~\ref{tab:extra}. As
baselines we report the majority-class proportion (0.504 on PTB-XL, 0.500 on MIT-BIH)
and, on PTB-XL, the ARI of uniformly random labels averaged over ten seeds (0.000).

\paragraph{Supervised ceiling.} The supervised ceiling is the accuracy, averaged over
five folds, of a logistic-regression probe with balanced class weights trained on the
same conditioned features, with folds stratified by class and grouped by patient
(PTB-XL) or record (MIT-BIH). The conditioner is fitted on all records before the folds
are formed. The ceiling lies in $[0,1]$, with the majority baseline as its chance level.
A high ceiling with a low ARI identifies a geometric rather than an
informational failure. To check that a low ceiling is not caused by the linearity of
the probe, a random-forest probe (200 trees, balanced class weights, same folds) was run
on the raw-flatten, hand-crafted, HuBERT-ECG and ECGFounder representations. Its accuracy
is reported for raw flattening, whose linear ceiling is closest to the majority baseline
and is the case the check addresses (Section~\ref{sec:sep}). The other three accuracies
were printed by the analysis script and were not recorded with the results.

\paragraph{Seeds, uncertainty and tests.} The $k$-means, spectral and GMM runs are
repeated over ten seeds (0 to 9), and each cell reports the mean. The seed changes only
the initialisation of these clusterers: the cohort, the MIT-BIH subsample, the
conv-AE (trained once) and the foundation embeddings are fixed. Ward agglomerative
clustering and HDBSCAN are deterministic and are run once. For every ten-seed cell we
report the seed standard deviation (SD) and a 95\% interval,
$\bar{x} \pm t_{0.975,9}\,\mathrm{SD}/\sqrt{10}$ with $t_{0.975,9}=2.262$
(Table~\ref{tab:seedci}). These intervals describe the variation due to clusterer
initialisation on a fixed cohort. They do not describe how the ARI would vary with a
different sample of patients. Four contrasts of the $k$-means ARI are tested with a
two-sided paired Wilcoxon signed-rank test across the ten seeds. The reported $p$-values
are exact sign-flip values: under the null hypothesis of the test each of the $2^{10}$
assignments of signs to the ten observed absolute paired differences is equally likely,
and tied absolute differences receive the mean of the ranks they span. When all ten
differences have the same sign, the observed assignment and its mirror image are the only
two that reach the most extreme value of the statistic, so the exact two-sided $p$-value is
$2/2^{10}=1/512\approx 0.002$, with or without ties. This is the smallest value the exact
test can give with ten pairs. The effect sizes are the difference of seed means $\Delta$,
the matched-pairs rank-biserial correlation $r=(T_+-T_-)/(m(m+1)/2)$, where $T_+$ and
$T_-$ are the sums of the ranks of the positive and the negative paired differences and
$m=10$ is the number of pairs, and Cliff's delta, the probability that a seed of the first
representation exceeds a seed of the second minus the probability of the reverse. Both $r$
and Cliff's delta range from $-1$ to 1.

\paragraph{Conditioners and stressors.} Because heavy-tailed features can dominate
Euclidean distance, every representation passes through a feature conditioner before
clustering. We sweep four conditioners, each fitted on the full representation matrix:
Yeo-Johnson power scaling, standard scaling, quantile scaling to a normal output
distribution, and robust scaling, and report power scaling as the default unless stated
otherwise. To check that fixing the cluster count to the number of classes does not
flatter the comparison, we also report adjusted mutual information, which is comparable
across cluster counts, for $K$ from 2 to 8, from one $k$-means run per $K$. Because the
natural SR-to-AF prevalence is imbalanced, we additionally evaluate on one 9:1 SR:AF
subsample: all SR records are kept, and $\lfloor 1282/9 \rfloor = 142$ AF records are
drawn at random with a fixed seed. The power conditioner is refitted on the subsample, and
one $k$-means run is made per representation.

\subsection{Federated robustness}

To probe deployment under data siloing we use standard federated $k$-means with
centroid aggregation in the style of federated averaging (FedAvg) \cite{fedavg2017}, and
we do not propose a new federated method. The 2\,542 records are partitioned, patient-grouped,
across ten simulated sites drawn from the single PTB-XL cohort, either independently and
identically distributed (IID), with each patient assigned to a site uniformly at random,
or with class proportions drawn from a Dirichlet distribution with concentration 0.3,
separately per class, so that sites differ sharply in arrhythmia prevalence (non-IID).
Sites with fewer than two records are dropped. Global centroids are initialised with $K=2$
records drawn at random from the pooled matrix. In each of 15 rounds every site starts
from the global centroids, runs three local Lloyd iterations, and returns its centroids,
which the server averages weighted by the number of the site's records assigned to each
centroid. Each record is finally assigned to its nearest global centroid. The power
conditioner is fitted on all records before partitioning. Each of the ten seeds draws a
new partition and a new initialisation. We report the global ARI relative to centralised
$k$-means of the same representation (ten initialisations per run, ten seeds), for the
raw-flatten, hand-crafted, HuBERT-ECG and ECGFounder representations.

\section{Results}\label{sec:results}

\subsection{PTB-XL benchmark: representation and clustering algorithm}

Table~\ref{tab:main} and Fig.~\ref{fig:heatmap} summarise the PTB-XL benchmark.
ECGFounder clusters best (ARI up to 0.885 with $k$-means), the hand-crafted features
are competitive (0.711 with Ward agglomerative clustering), HuBERT-ECG is weak (0.195),
and raw flattening and the conv-AE are indistinguishable from a majority-class assignment
(ARI at most 0.013). On this task the spread across clustering algorithms within a strong
representation is smaller than the spread across representations, so the representation
is the deciding factor here. Section~\ref{sec:sep} discusses why the representations
differ.

\begin{table}[t]
\centering
\caption{PTB-XL SR-vs-AF ($n = 2\,542$ records: 1\,282 SR, 1\,260 AF), power scaling:
clustering ARI and the supervised ceiling (patient-grouped five-fold accuracy of a
logistic-regression probe with balanced class weights). $k$-means, spectral and GMM cells
are means over ten seeds. Agglomerative (agglom., Ward linkage) and HDBSCAN cells are single
runs of deterministic algorithms. Numbers in parentheses are dimensionalities. Sup.\
ceiling: supervised ceiling. Majority-class baseline 0.504, ARI of random labels 0.000.
Seed SDs are at most 0.0013 except raw-flatten GMM (0.0133) and ECGFounder GMM (0.0116).
SDs and 95\% intervals of all ten-seed cells are in Table~\ref{tab:seedci}. The four
tested $k$-means contrasts have exact two-sided $p=1/512\approx 0.002$ (paired Wilcoxon
signed-rank test, the smallest value the exact test can give with ten seeds), with
rank-biserial $r=1$ and Cliff's delta 1.
The ordering holds for $K=2$ to $8$ under adjusted mutual information.}
\label{tab:main}
\begin{tabular}{lcccccc}
\toprule
Representation & $k$-means & agglom. & spectral & GMM & HDBSCAN & sup.\ ceiling \\
\midrule
raw-flatten (1\,992)    & 0.000 & 0.000 & 0.000 & 0.013 & 0.000 & 0.519 \\
conv-AE (32)            & 0.004 & 0.008 & 0.001 & 0.005 & 0.000 & 0.678 \\
hand-crafted (16)       & 0.644 & \textbf{0.711} & 0.672 & 0.544 & 0.060 & 0.942 \\
HuBERT-ECG (768)        & 0.174 & 0.134 & 0.016 & 0.195 & 0.000 & 0.899 \\
ECGFounder (1\,024)     & \textbf{0.885} & 0.709 & 0.860 & 0.851 & 0.137 & 0.983 \\
\bottomrule
\end{tabular}
\end{table}

\begin{figure}[t]
\centering
\includegraphics[width=0.78\linewidth]{figures/fig1_ari_heatmap.pdf}
\caption{PTB-XL SR-vs-AF clustering ARI by representation and clusterer (power scaling,
$n = 2\,542$ records, 1\,282 SR and 1\,260 AF). Cell values are the Table~\ref{tab:main}
values rounded to two decimals. Axis labels: kmean, $k$-means; agglo, Ward agglomerative;
spect, spectral; gmm, GMM; hdbsc, HDBSCAN.}
\label{fig:heatmap}
\end{figure}

Across the ten seeds the $k$-means, spectral and GMM solutions vary little: the seed SD
of the ARI is at most 0.0013 in thirteen of the fifteen ten-seed cells, and 0.0133 and
0.0116 for the raw-flatten and ECGFounder GMM cells (Table~\ref{tab:seedci}). For example,
the 95\% interval of the ECGFounder $k$-means ARI is [0.885, 0.886] and that of the
hand-crafted $k$-means ARI is the single value 0.644, because all ten seeds give the same
partition. The gaps between representations are therefore stable over clusterer
initialisation. These intervals do not cover the sampling of patients. A paired Wilcoxon
signed-rank test on the $k$-means seeds was applied to four contrasts: ECGFounder over
hand-crafted ($\Delta = 0.242$), hand-crafted over HuBERT-ECG ($\Delta = 0.470$),
hand-crafted over raw flattening ($\Delta = 0.644$), and HuBERT-ECG over raw flattening
($\Delta = 0.174$), where $\Delta$ is the difference of the unrounded seed means. In each
contrast all ten paired differences are positive, so the exact two-sided $p$-value is
$1/512\approx 0.002$ (Section~\ref{sec:clust}) and the matched-pairs rank-biserial
correlation is $r=1$. Every seed of the first representation exceeds every seed of the
second, so Cliff's delta is 1. The $p$-value and both effect sizes are at their extreme
values and carry the same information. The informative effect size is the ARI difference
$\Delta$ itself. The ordering does not
depend on fixing the cluster count to the number of classes: with adjusted mutual
information for $K$ from 2 to 8 (one $k$-means run per $K$), ECGFounder leads at every $K$
(0.83 falling to 0.45), the hand-crafted features follow (0.56 to 0.33), and HuBERT-ECG
and raw flattening stay low (below 0.14 and near 0).

The two foundation models do not behave as one category. ECGFounder clusters strongly while
HuBERT-ECG does not, which matches the report of \cite{lookingbeyond2026} that some ECG
foundation embeddings organise by clinical label and others by dataset. We therefore name
each model and do not generalise to foundation models as a class. HDBSCAN gives the
lowest or joint-lowest ARI of the five clusterers for every PTB-XL representation (at most
0.137). Under power scaling it labels every record as noise for raw flattening, the
conv-AE and HuBERT-ECG (no cluster found, ARI 0.000). For the hand-crafted features and
ECGFounder it finds two clusters, and the stored purity and cluster sizes place the noise
fraction between 0.720 and 0.833 (hand-crafted) and between 0.595 and 0.774 (ECGFounder).
Because the noise points enter the ARI as one cluster, these HDBSCAN scores describe a
partition in which most records form a single noise group. Spectral clustering collapses
where the nearest-neighbour affinity graph is poorly connected. We report both rather than
tune them away.

\subsection{Separability is not clusterability}\label{sec:sep}

The supervised ceilings (last column of Table~\ref{tab:main}, Fig.~\ref{fig:sep})
expose the gap. Raw flattening carries almost no linearly recoverable label information
(ceiling 0.519, barely above the 0.504 majority), consistent with its chance-level
clustering. The conv-AE is more separable (0.678) yet still clusters at chance (ARI at
most 0.008, Ward agglomerative): a properly trained reconstruction autoencoder encodes
the class without organising it into clusters. HuBERT-ECG shows the largest gap between
separability and clusterability: a linear probe reaches an accuracy of 0.899, yet the best
clustering reaches an ARI of 0.195 (GMM). Accuracy (chance level 0.504 here) and ARI
(chance level 0) are on different scales, so the two are compared ordinally. A high
ceiling with a low ARI means that the label is present in the embedding and can be
recovered when labels are available for training, whereas label-free clustering of the
embedding does not recover it. ECGFounder and the hand-crafted features achieve both high
separability and high clusterability. The discriminative information exists in the weaker
embeddings, but it does not lie along their dominant axes of variation, the
separability-versus-clusterability effect reported for self-supervised representations \cite{sslclustering2025}. This is not an
artefact of a linear probe: a random forest on the raw-flatten representation reaches
accuracy 0.71, well above its linear ceiling of 0.519 and far above its clustering ARI (at
most 0.013), so even a non-linear decoder recovers a label that clustering does not.

The representations differ for reasons that the ceilings and the construction of each
representation make visible. The hand-crafted features are competitive because SR and AF
differ by definition in the regularity of the rhythm, and eight of the 16 features are
computed from the detected R peaks. The high ceiling (0.942) and ARI are consistent with
the label lying along a few dominant feature axes. Raw flattening matches the majority
assignment because the flattened strips are not aligned in time across recordings, so
Euclidean distance between them reflects where the beats fall in the strip and how large
they are rather than the rhythm. Its linear ceiling (0.519) is barely above the majority
baseline, and its Ward agglomerative solution places all records but one in one cluster
(Section~\ref{sec:trap}). The conv-AE minimises reconstruction error, which rewards
reproducing signal energy and waveform shape. Its code carries the label weakly (ceiling
0.678) and does not organise it into clusters. ECGFounder has the highest ceiling (0.983),
and the four clusterers other than HDBSCAN recover its label (ARI 0.709 to 0.885), which is
consistent with its label information lying along its dominant variation. A plausible reason is that ECGFounder was trained on expert diagnostic labels
\cite{ecgfounder2024}, so that its features are organised by diagnostic categories. For
HuBERT-ECG the label is present (ceiling 0.899) but the dominant directions partitioned by
$k$-means and GMM are not the label direction. We offer three hypotheses, none tested
here. First, its self-supervised pretraining objective does not use diagnostic labels, so
its dominant directions need not follow rhythm class, consistent with the report that some
foundation embeddings organise by dataset rather than by label \cite{lookingbeyond2026}.
Second, in this study it receives only the first 5\,s of each record, half of the strip
seen by the other representations, which leaves fewer beats from which rhythm
irregularity can be read. Third, its input here is neither band-pass filtered nor scaled
to $[-1,1]$ per lead, as its authors' preprocessing does (Section~\ref{sec:repr}), so it
does not receive amplitudes on the scale used in its pretraining.

\begin{figure}[t]
\centering
\includegraphics[width=0.60\linewidth]{figures/fig3_sep_vs_clust.pdf}
\caption{Separability is not clusterability ($n = 2\,542$ PTB-XL records). Each point is
one representation: the horizontal axis is the supervised ceiling (patient-grouped
five-fold accuracy of the linear probe, ``acc'' in the axis label), the vertical axis the
highest ARI over the five clusterers in Table~\ref{tab:main}. The plotted conv-AE and
hand-crafted values are single runs of the deterministic Ward agglomerative clustering,
and the seed SDs of the other three are at most 0.0133 (Table~\ref{tab:seedci}). The
dashed line labelled ``parity'' joins (0.5, 0.5) and (1, 1). Because accuracy has chance
level 0.504 and ARI has chance level 0, it is a visual guide only, and the vertical distance of a point below it is
not a calibrated measure. Points far below the line are separable but not clusterable.}
\label{fig:sep}
\end{figure}

\subsection{The silhouette trap}\label{sec:trap}

Internal indices mislead here in a specific, reproducible way. Raw flattening with
agglomerative clustering attains a silhouette of 0.81 while all records but one
fall into one cluster (largest-cluster fraction 1.00) and the ARI is 0.000. The
hand-crafted representation, which actually recovers the classes (ARI 0.711, balanced
clusters), has a \emph{lower} silhouette of 0.28 (Fig.~\ref{fig:trap}). A degenerate
one-cluster solution maximises compactness while carrying no class information. This is
the artefact behind the original project's silhouette of 0.89, and it is why
chance-corrected external metrics, read together with cluster-size sanity checks, are
necessary.

\begin{figure}[t]
\centering
\includegraphics[width=0.62\linewidth]{figures/fig2_silhouette_trap.pdf}
\caption{The silhouette trap ($n = 2\,542$ PTB-XL records, power scaling). Silhouette
(horizontal axis, computed on non-noise points) against chance-corrected ARI (vertical
axis). Marker size is the largest-cluster fraction. Each point is one clustering run of
one representation: ten seeds each for $k$-means, spectral and GMM and one run each for
Ward agglomerative and HDBSCAN, 160 runs in total, of which 157 are drawn (the three
HDBSCAN runs that labelled every record as noise have no silhouette). Runs that give the
same partition coincide. The degenerate one-cluster solutions (large markers, bottom
right) attain high silhouette at zero ARI.}
\label{fig:trap}
\end{figure}

\subsection{Conditioning matters}

The feature conditioner changes the outcome materially. For the hand-crafted features the
mean ARI over all 32 runs per conditioner (ten seeds each of $k$-means, spectral and GMM,
one run each of Ward agglomerative and HDBSCAN) is 0.63 under standard scaling and 0.61
under power scaling, but
falls to 0.36 under robust scaling and 0.23 under quantile scaling. The effect is not
confined to hand-crafted features: ECGFounder's mean ARI drops from 0.87 (standard) to
0.09 under robust scaling. Robust scaling, a common default for physiological signals,
leaves heavy arrhythmic tails unbounded so that a few outliers dominate Euclidean
distance and the partition collapses. Conditioning should be reported and swept, not
assumed.

\subsection{Generalisation to MIT-BIH beat morphology}

The ordering of the representations by their best ARI replicates on a different dataset
and a different task type (Table~\ref{tab:mitbih}). On the N-vs-V task the 21
beat-morphology features cluster well (ARI 0.783, Ward agglomerative). The raw beat window
reaches an ARI of at most 0.049 with $k$-means, Ward agglomerative clustering and GMM, and
0.107 with HDBSCAN, despite beats being R-peak aligned. On MIT-BIH, unlike on PTB-XL,
HDBSCAN therefore gives the highest ARI for the raw beat window, which is still far below
the 0.783 of the beat-morphology features. Spectral clustering places almost all beats in
one cluster for all three representations (ARI at most 0.002, largest-cluster fraction at
least 0.978, Table~\ref{tab:extra}). The conv-AE is again the clean
separability-versus-clusterability case: a supervised ceiling of 0.927 but an ARI of
0.110 with $k$-means and at most 0.127 over the five clusterers. That the same pattern appears for rhythm on
PTB-XL and for morphology on MIT-BIH supports the representation-level reading for these
two tasks.

\begin{table}[t]
\centering
\footnotesize
\caption{MIT-BIH N-vs-V beat morphology ($n = 3\,000$ beats, 1\,500 per class, from 47
records), power scaling: clustering ARI. $k$-means, spectral and GMM cells are means over
ten seeds. Ward agglomerative and HDBSCAN cells are single runs of deterministic
algorithms. Numbers in parentheses are dimensionalities. Supervised ceilings: beat-morphology 0.925, conv-AE
0.927, raw-beat 0.741. Majority-class baseline 0.500. The foundation models were not
applied (Section~\ref{sec:repr}).}
\label{tab:mitbih}
\begin{tabular}{lccc}
\toprule
Clusterer & beat-morphology (21) & conv-AE (32) & raw-beat (396) \\
\midrule
$k$-means     & 0.750 & 0.110 & 0.018 \\
Ward agglomerative & \textbf{0.783} & 0.071 & 0.049 \\
spectral      & 0.001 & 0.002 & 0.000 \\
GMM           & 0.679 & 0.127 & 0.018 \\
HDBSCAN       & 0.074 & 0.036 & 0.107 \\
\bottomrule
\end{tabular}
\end{table}

\subsection{Robustness under deployment stress: imbalance and federation}

Two simulated stressors closer to deployment change the comparison between ECGFounder
and the hand-crafted features. First, class imbalance. On one 9:1 SR:AF subsample
(1\,282 SR and 142 AF records, majority 0.90), every representation's ARI falls, because
isolating a 10\% minority cluster with a partitional method is intrinsically hard, but the
hand-crafted features retain the most structure (ARI 0.232, AMI 0.214) and
\emph{overtake} ECGFounder (ARI 0.060, AMI 0.143), with HuBERT-ECG and raw flattening near
zero. These are single-run values, from one $k$-means run on one subsample. ECGFounder's balanced-data lead does not hold on this subsample.

Second, federation. Under IID federation across ten simulated sites, every representation
retains essentially its centralised ARI, which validates the federated procedure
(Table~\ref{tab:fed}, Fig.~\ref{fig:fed}). Non-IID partitions lower the mean ARI of
ECGFounder, the hand-crafted features and HuBERT-ECG, but unequally. ECGFounder keeps the
higher mean ARI (0.612 against 0.591 for the hand-crafted features), and its centralised
advantage shrinks from 0.242 to 0.021. The hand-crafted features keep a larger share of
their centralised ARI (retention 0.918 against 0.692) and vary far less across the ten
seeds, each of which draws a new site partition and initialisation: the seed SD is 0.071
against 0.369, and the 95\% intervals over seeds are [0.541, 0.642] for the hand-crafted
features and [0.348, 0.876] for ECGFounder (Table~\ref{tab:seedci}). The two intervals
overlap, so these intervals do not separate the two representations on mean ARI.
ECGFounder is ahead on the mean, and the hand-crafted features are markedly more stable
across partitions. These results concern ten sites simulated from one dataset by Dirichlet
sampling of class proportions.

\begin{table}[t]
\centering
\caption{Federated $k$-means on PTB-XL ($n = 2\,542$ records, ten simulated sites):
mean ARI over ten seeds. Each seed draws a new site partition and initialisation.
Retention is the ratio of the non-IID to the centralised mean ARI, computed from unrounded
means. N/A: not interpreted, because both raw-flatten ARIs are below 0.001 (the ratio of
the stored means is 1.505). Seed SDs and 95\% intervals are in Table~\ref{tab:seedci}.}
\label{tab:fed}
\begin{tabular}{lcccc}
\toprule
Representation & centralised & federated IID & federated non-IID & retention \\
\midrule
ECGFounder   & 0.885 & 0.885 & 0.612 & 0.692 \\
hand-crafted & 0.644 & 0.643 & 0.591 & \textbf{0.918} \\
HuBERT-ECG   & 0.174 & 0.172 & 0.106 & 0.611 \\
raw-flatten  & 0.000 & 0.000 & 0.000 & N/A \\
\bottomrule
\end{tabular}
\end{table}

\begin{figure}[t]
\centering
\includegraphics[width=0.78\linewidth]{figures/fig4_federated.pdf}
\caption{Federated $k$-means across ten simulated sites (labelled ``clients'' in the plot
title, $n = 2\,542$ PTB-XL records): mean ARI over ten seeds, without error bars (seed
SDs and intervals in Table~\ref{tab:seedci}). The raw-flatten bars have height 0.000.
Under non-IID partitions ECGFounder keeps the higher mean ARI, and the hand-crafted
features keep the larger share of their centralised ARI.}
\label{fig:fed}
\end{figure}

\section{Discussion}\label{sec:discussion}

The observations of this study are the values in Tables~\ref{tab:main} to~\ref{tab:fed},
obtained on two binary tasks, with five representations, two foundation models and a
simulated federation. The practical reading below is a decision rule that follows from
these observations and is a recommendation conditional on these settings. If a strong public foundation model can be run on balanced data
in the target 12-lead format, ECGFounder gave the best unsupervised separation here. If
the data are imbalanced, or if interpretability or low dimensionality matters, the 16
hand-crafted features were the better choice in these experiments. Their best ARI is 0.174
below ECGFounder's on balanced data (0.711 against 0.885), and they overtake it on the 9:1
subsample (0.232 against 0.060). If the data are siloed across sites, the choice depends on
whether the mean or the stability across partitions matters more. Under simulated non-IID
federation ECGFounder has the higher mean ARI (0.612 against 0.591), and the hand-crafted
features vary far less across site partitions (seed SD 0.071 against 0.369).

The input format also matters. In this study the foundation models received fixed-length
12-lead strips (5\,s for HuBERT-ECG, 10\,s for ECGFounder) and were not applied to the
550\,ms two-channel MIT-BIH beats. A beat-level hand-crafted feature set (the 21 MIT-BIH
features, which differ from the 16 PTB-XL features) could be computed for that format. The
16 PTB-XL features are computed from lead II alone (Section~\ref{sec:repr}), so they can be
computed wherever lead II is recorded, but their comparison with the foundation models on
such input was not evaluated. Whether the foundation models can embed shorter or
fewer-lead inputs was not tested here. The ECGFounder authors describe an extension of
their model to single-lead ECGs \cite{ecgfounder2024}. The raw-flatten and
reconstruction-autoencoder representations, including the original project's design,
reached ARIs of at most 0.013 on PTB-XL and at most 0.127 on MIT-BIH with the
configurations tested here. On PTB-XL they are indistinguishable from predicting the
majority class.

Deployment also depends on cost. Table~\ref{tab:cost} lists, for each PTB-XL
representation, its dimensionality, whether it is trained, frozen or computed without
training, the time per record and the peak inference memory. \textcolor{red}{[AUTHOR: one or two
sentences relating the measured costs to the decision rule above, once the cost columns of
Table~\ref{tab:cost} are filled.]}

\begin{table}[t]
\centering
\footnotesize
\caption{Cost of each PTB-XL representation. Dimensionality and training status follow
from Table~\ref{tab:repr}. Trained: trained on the cohort. Time/record: time per record
for feature extraction (raw flattening, hand-crafted) or embedding (conv-AE, HuBERT-ECG,
ECGFounder, including resampling). Peak memory: peak inference memory.
\textcolor{red}{[AUTHOR: hardware, batch size and numeric precision of the timing runs.]}}
\label{tab:cost}
\begin{tabular}{lrlll}
\toprule
Representation & Dimensions & Training & Time/record & Peak memory \\
\midrule
raw-flatten  & 1\,992 & none   & \textcolor{red}{[AUTHOR]} & \textcolor{red}{[AUTHOR]} \\
conv-AE      & 32     & trained & \textcolor{red}{[AUTHOR]} & \textcolor{red}{[AUTHOR]} \\
hand-crafted & 16     & none   & \textcolor{red}{[AUTHOR]} & \textcolor{red}{[AUTHOR]} \\
HuBERT-ECG   & 768    & frozen & \textcolor{red}{[AUTHOR]} & \textcolor{red}{[AUTHOR]} \\
ECGFounder   & 1\,024 & frozen & \textcolor{red}{[AUTHOR]} & \textcolor{red}{[AUTHOR]} \\
\bottomrule
\end{tabular}
\end{table}

The result also reframes the original master's project. Its conclusion that the
bottleneck is the representation rather than the clustering algorithm is correct, and we
have quantified it. Its headline metrics were the artefacts of two avoidable choices,
purity without a baseline and silhouette on a collapsed embedding. Both are corrected
here, and both recur often enough in applied clustering to be worth stating plainly.

The contribution is methodological at the level of evaluation: a protocol for judging
unsupervised ECG representations that combines chance-corrected external indices with
majority and random baselines, a supervised ceiling per representation computed under the
same patient grouping, cluster-size checks, a conditioner sweep, and tests under two
simulated deployment stressors. It introduces no new representation, clustering algorithm
or federated method. The separability-versus-clusterability phenomenon is
known \cite{sslclustering2025}, the federated procedure is standard \cite{fedavg2017},
and the representations are existing tools. Relative to the closest prior benchmark
\cite{lookingbeyond2026}, the protocol adds the hand-crafted baseline and the MIT-BIH
dataset, and it applies the same controls to non-foundation and foundation
representations alike.

The foundation-model axis is bounded to two named models. Both have publicly released
weights, are used here as frozen feature extractors, and accept 12-lead input in the
format of PTB-XL. They differ in pretraining, HuBERT-ECG being self-supervised and
ECGFounder trained on expert diagnostic labels \cite{hubertecg2024,ecgfounder2024}.
\textcolor{red}{[AUTHOR: the reasons these two models were selected over other publicly
available ECG foundation models.]} The bound affects the conclusions in three ways. Every
foundation-model statement in this paper concerns these two models only. The comparison
between hand-crafted features and a foundation model is a comparison with ECGFounder. A
foundation model that clusters better than ECGFounder under imbalance or under simulated
federation would change the decision rule above. Other foundation models, including some
reported to cluster well \cite{lookingbeyond2026}, were not evaluated, and our claims are
scoped accordingly.

\section{Limitations}\label{sec:limitations}

ECGFounder expects 500\,Hz input, and because bulk streaming of the 500\,Hz records was
rate-limited at the source, the main benchmark embeds the 100\,Hz cohort FFT-resampled to
the expected length of 5\,000 samples per lead. A consistency check on 300 records (the
first 150 SR and the first 150 AF records in the order in which the cohort was drawn), one
$k$-means run each, compared native 500\,Hz input with upsampled 100\,Hz input: the ARI
was 0.896 at 500\,Hz and 0.909 after upsampling, a difference of 0.013 in favour of the
upsampled input. The check covers one subset and one run, so it bounds the effect of
upsampling only for that subset. Neither foundation model received the filtering its
authors apply (Section~\ref{sec:repr}). HuBERT-ECG received unfiltered values that were not
scaled to $[-1,1]$, in place of values band-pass filtered from 0.05 to 47\,Hz and scaled
per lead, and ECGFounder received unfiltered values in place of high-pass, low-pass and
notch-filtered ones. Whether matching these steps changes the ARI of either model was not
tested. We evaluate two foundation models. Other foundation models, including some
reported to cluster well \cite{lookingbeyond2026}, are not included, which is why we name
the two we test rather than generalising. The study evaluates two binary tasks, SR versus
AF on PTB-XL and N versus V on MIT-BIH, and the results apply to these two tasks. The
positive PTB-XL class includes atrial flutter as well as
atrial fibrillation. The hand-crafted PTB-XL features use lead II only, and other leads
were not evaluated for them. The MIT-BIH arm omits the foundation models because their
input in this study is a strip of 5\,s or 10\,s, much longer than a 550\,ms beat window.
The federated study is a simulation on a single physical dataset, with sites formed by
random or Dirichlet-sampled partitions of one cohort. It does not model differences in
devices, acquisition protocols or populations between institutions, and real
multi-institution data were not evaluated. The simulation also initialises the global
centroids from records drawn from the pooled matrix and fits the power conditioner on all
records before partitioning, neither of which a real federation could do without sharing
data between sites. Finally, the reported intervals describe initialisation and partition
variability on the cohorts used.

\section{Conclusion}\label{sec:conclusion}

For the two binary tasks evaluated here, SR versus AF on PTB-XL and N versus V on
MIT-BIH, the representation decided whether clustering worked more than the clustering
algorithm did, and chance-corrected external evaluation was required to see it. On
balanced data ECGFounder leads. On one 9:1 class-imbalanced subsample the 16
interpretable hand-crafted features overtake ECGFounder. Under simulated non-IID
federation ECGFounder
keeps the higher mean ARI, while the hand-crafted features retain a larger share of their
centralised ARI and vary far less across site partitions. Raw-flatten and
reconstruction-autoencoder representations stay near chance on PTB-XL regardless of the
clustering method. Internal indices and feature conditioning are not neutral and must be
reported. In these experiments the deciding choices were the representation and the
evaluation protocol rather than the clustering algorithm.

\section*{Code and data availability}
PTB-XL (version 1.0.3) \cite{ptbxl2020,ptbxl2022physionet} and the MIT-BIH Arrhythmia Database (version 1.0.0) \cite{mitbih2001,mitbih2005physionet} are public
on PhysioNet. The benchmark code, with the exact cohort
definitions, is available at \url{https://github.com/TheAxeC/ecg-representation-clustering}
(commit 78472d3). \textcolor{red}{[AUTHOR: tagged release or archived DOI of the code
version that produced the reported results.]} The pretrained weights are those of
HuBERT-ECG at \url{https://huggingface.co/Edoardo-BS/hubert-ecg-base} and of ECGFounder
at \url{https://huggingface.co/PKUDigitalHealth/ECGFounder} (file
\texttt{12\_lead\_ECGFounder.pth}). The preprocessing and model-definition code of the
two models' authors referred to in Section~\ref{sec:repr} is at
\url{https://github.com/Edoar-do/HuBERT-ECG} (file \texttt{hubert\_ecg/utils.py},
function \texttt{ecg\_preprocessing}, and \texttt{hubert\_ecg/dataset.py}) and
\url{https://github.com/PKUDigitalHealth/ECGFounder} (files \texttt{dataset.py} and
\texttt{finetune\_model.py}). Implementation settings are listed in Table~\ref{tab:impl}.

\section*{Ethics}
This study uses only de-identified, publicly available databases and required no new
human-subjects data collection.

\section*{CRediT author contributions}
\textbf{Dongho Chun}: methodology, software, investigation, writing -- original draft.
\textbf{Axel Faes}: conceptualization, methodology, supervision, writing -- review and
editing.

\section*{Funding}
This work was supported by the Flanders AI Research Program (FAIR), an Impulse Program of the Flemish Government (Flemish Ministry of Economy, Science and Innovation).

\section*{Declaration of competing interest}
The authors declare that they have no known competing financial interests or personal relationships that could have appeared to influence the work reported in this paper.

\section*{Declaration of generative AI and AI-assisted technologies in the writing process}
During the preparation of this work the authors used a generative AI assistant (Claude) to support software implementation and to edit and improve the language of the manuscript. After using this tool, the authors reviewed and edited the content as needed and take full responsibility for the content of the publication.

\appendix
\setcounter{table}{0}

\section{Implementation settings and supplementary results}\label{app:impl}

Table~\ref{tab:impl} lists the settings of every step. Settings not listed are the
defaults of the library named. Table~\ref{tab:seedci} gives the seed SD and the 95\%
interval of every ten-seed cell of Tables~\ref{tab:main}, \ref{tab:mitbih} and
\ref{tab:fed}. Table~\ref{tab:extra} gives the further external indices of every cell of
Tables~\ref{tab:main} and~\ref{tab:mitbih}.

\begin{table}[ht]
\centering
\footnotesize
\caption{Implementation settings. $K=2$ throughout. Seeds are 0 to 9 unless a fixed seed
is stated.}
\label{tab:impl}
\begin{tabular}{p{0.22\linewidth}p{0.70\linewidth}}
\toprule
Step & Setting \\
\midrule
$k$-means & scikit-learn \texttt{KMeans}, 10 initialisations per run, seed as random state \\
Ward agglomerative & \texttt{AgglomerativeClustering}, Ward linkage; deterministic, one run \\
Spectral & \texttt{SpectralClustering}, nearest-neighbour affinity, $k$-means label
assignment, seed as random state \\
GMM & \texttt{GaussianMixture}, diagonal covariance, covariance regularisation
$10^{-4}$, seed as random state \\
HDBSCAN & \texttt{sklearn.cluster.HDBSCAN}, minimum cluster size $\max(10,\lfloor
n/50\rfloor)$; deterministic, one run; noise label kept as one cluster in the external
indices \\
Conditioners & power: \texttt{PowerTransformer} (Yeo-Johnson); standard:
\texttt{StandardScaler}; quantile: \texttt{QuantileTransformer} with normal output;
robust: \texttt{RobustScaler}; each fitted on the full matrix \\
Linear probe & \texttt{LogisticRegression}, at most 2\,000 iterations, balanced class
weights; \texttt{StratifiedGroupKFold}, 5 folds, no shuffling; accuracy \\
Random-forest probe & 200 trees, balanced class weights, fixed seed 0; same folds;
accuracy \\
Conv-AE encoder & three one-dimensional convolutions (32, 64, 128 channels; kernel 7, stride 2,
padding 3) with LeakyReLU, adaptive average pooling, linear layer to 32 units \\
Conv-AE decoder & linear layer from 32 units to $128 \times \lfloor T/8 \rfloor$ ($T=2\,500$),
three transposed one-dimensional convolutions (64, 32 and input channels; kernel 8, stride 2,
padding 3) with LeakyReLU, output cropped or padded to $T$ \\
Conv-AE training & Adam, learning rate $10^{-3}$, mean squared error, batch 64, at most
50 epochs (PTB-XL) or 40 (MIT-BIH); random 10\% validation split; early stopping after 8
epochs without a validation-loss improvement of $10^{-5}$, best weights restored; trained
once with seed 0 on all records, and the embedding is computed for all records \\
NeuroKit2 & \texttt{ecg\_clean} and \texttt{ecg\_peaks} with default methods,
\texttt{ecg\_delineate} with the discrete wavelet method, at 100\,Hz on lead II \\
HuBERT-ECG & checkpoint \texttt{Edoardo-BS/hubert-ecg-base} loaded into a
\texttt{transformers} \texttt{HubertModel} without the pretraining heads; mean over time
of the last hidden state, no attention mask; batch 64; 32-bit floating point; no gradient \\
ECGFounder & checkpoint \texttt{PKUDigitalHealth/ECGFounder}, \texttt{Net1D} backbone
(filters 64, 160, 160, 400, 400, 1\,024, 1\,024; blocks 2, 2, 2, 3, 3, 4, 4; kernel 16;
stride 2; group width 16), built with batch normalisation and dropout disabled as in the
authors' model definition, classification layer removed; time average of the last stage; 32-bit floating point; no
gradient; batch size \textcolor{red}{[AUTHOR]} \\
Checkpoint revisions & \textcolor{red}{[AUTHOR: Hugging Face commit of each checkpoint]} \\
Federation & 10 sites; Dirichlet concentration 0.3 per class (non-IID) or uniform
assignment of patients (IID); 15 rounds of 3 local Lloyd iterations; centroids averaged
weighted by assigned counts; sites with fewer than 2 records dropped; per-site sizes
\textcolor{red}{[AUTHOR]} \\
9:1 subsample & all 1\,282 SR records and 142 AF records drawn with fixed seed 0; one
$k$-means run (10 initialisations, seed 0) per representation \\
Software & lower bounds: numpy 1.24, scipy 1.10, pandas 2.0, scikit-learn 1.3, wfdb 4.1,
neurokit2 0.2.7, torch 2.0, transformers 4.40; versions used
\textcolor{red}{[AUTHOR]} \\
\bottomrule
\end{tabular}
\end{table}

\begin{table}[ht]
\centering
\footnotesize
\caption{Seed variation of every ten-seed cell (power scaling): mean ARI, seed SD
(sample SD over ten seeds) and 95\% interval $\bar{x} \pm 2.262\,\mathrm{SD}/\sqrt{10}$.
For PTB-XL and MIT-BIH centralised cells the seed changes only the clusterer
initialisation. For federated cells it changes the site partition and the initialisation.
The intervals do not cover the sampling of patients or records. Values are rounded to
three decimals (SD to four).}
\label{tab:seedci}
\begin{tabular}{llrrc}
\toprule
Table and representation & Setting & Mean & SD & 95\% interval \\
\midrule
\multicolumn{5}{l}{PTB-XL (Table~\ref{tab:main})} \\
raw-flatten  & $k$-means & 0.000 & 0.0002 & [0.000, 0.000] \\
raw-flatten  & spectral  & 0.000 & 0.0000 & [0.000, 0.000] \\
raw-flatten  & GMM       & 0.013 & 0.0133 & [0.003, 0.022] \\
conv-AE      & $k$-means & 0.004 & 0.0000 & [0.004, 0.004] \\
conv-AE      & spectral  & 0.001 & 0.0000 & [0.001, 0.001] \\
conv-AE      & GMM       & 0.005 & 0.0001 & [0.005, 0.005] \\
hand-crafted & $k$-means & 0.644 & 0.0000 & [0.644, 0.644] \\
hand-crafted & spectral  & 0.672 & 0.0000 & [0.672, 0.672] \\
hand-crafted & GMM       & 0.544 & 0.0000 & [0.544, 0.544] \\
HuBERT-ECG   & $k$-means & 0.174 & 0.0000 & [0.174, 0.174] \\
HuBERT-ECG   & spectral  & 0.016 & 0.0000 & [0.016, 0.016] \\
HuBERT-ECG   & GMM       & 0.195 & 0.0002 & [0.195, 0.195] \\
ECGFounder   & $k$-means & 0.885 & 0.0012 & [0.885, 0.886] \\
ECGFounder   & spectral  & 0.860 & 0.0000 & [0.860, 0.860] \\
ECGFounder   & GMM       & 0.851 & 0.0116 & [0.843, 0.859] \\
\midrule
\multicolumn{5}{l}{MIT-BIH (Table~\ref{tab:mitbih})} \\
beat-morphology & $k$-means & 0.750 & 0.0610 & [0.706, 0.793] \\
beat-morphology & spectral  & 0.001 & 0.0005 & [0.000, 0.001] \\
beat-morphology & GMM       & 0.679 & 0.0006 & [0.678, 0.679] \\
conv-AE         & $k$-means & 0.110 & 0.0179 & [0.098, 0.123] \\
conv-AE         & spectral  & 0.002 & 0.0015 & [0.001, 0.003] \\
conv-AE         & GMM       & 0.127 & 0.1505 & [0.019, 0.234] \\
raw-beat        & $k$-means & 0.018 & 0.0000 & [0.018, 0.018] \\
raw-beat        & spectral  & 0.000 & 0.0003 & [0.000, 0.001] \\
raw-beat        & GMM       & 0.018 & 0.0000 & [0.018, 0.018] \\
\midrule
\multicolumn{5}{l}{Federated $k$-means on PTB-XL (Table~\ref{tab:fed})} \\
ECGFounder   & centralised & 0.885 & 0.0012 & [0.885, 0.886] \\
ECGFounder   & IID         & 0.885 & 0.0020 & [0.883, 0.886] \\
ECGFounder   & non-IID     & 0.612 & 0.3691 & [0.348, 0.876] \\
hand-crafted & centralised & 0.644 & 0.0000 & [0.644, 0.644] \\
hand-crafted & IID         & 0.643 & 0.0009 & [0.643, 0.644] \\
hand-crafted & non-IID     & 0.591 & 0.0705 & [0.541, 0.642] \\
HuBERT-ECG   & centralised & 0.174 & 0.0000 & [0.174, 0.174] \\
HuBERT-ECG   & IID         & 0.172 & 0.0019 & [0.170, 0.173] \\
HuBERT-ECG   & non-IID     & 0.106 & 0.0416 & [0.076, 0.136] \\
raw-flatten  & centralised & 0.000 & 0.0002 & [0.000, 0.000] \\
raw-flatten  & IID         & 0.000 & 0.0003 & [0.000, 0.000] \\
raw-flatten  & non-IID     & 0.000 & 0.0002 & [0.000, 0.000] \\
\bottomrule
\end{tabular}
\end{table}

\begin{table}[ht]
\centering
\footnotesize
\caption{Further external indices (power scaling): normalised mutual information (NMI),
Fowlkes-Mallows index (FMI), purity and largest-cluster fraction (LCF), each the mean
over ten seeds for $k$-means, spectral and GMM and a single run for Ward agglomerative and
HDBSCAN. HDBSCAN noise counts as one cluster in NMI and FMI, as misassigned in purity, and
is excluded from the cluster sizes in LCF. n/a: no non-noise cluster.}
\label{tab:extra}
\begin{tabular}{llcccc}
\toprule
Representation & Clusterer & NMI & FMI & purity & LCF \\
\midrule
\multicolumn{6}{l}{PTB-XL ($n = 2\,542$, majority 0.504)} \\
raw-flatten  & $k$-means & 0.001 & 0.685 & 0.506 & 0.967 \\
raw-flatten  & agglom.   & 0.001 & 0.707 & 0.505 & 1.000 \\
raw-flatten  & spectral  & 0.001 & 0.703 & 0.506 & 0.995 \\
raw-flatten  & GMM       & 0.012 & 0.633 & 0.542 & 0.862 \\
raw-flatten  & HDBSCAN   & 0.000 & 0.707 & 0.000 & n/a \\
conv-AE      & $k$-means & 0.003 & 0.503 & 0.532 & 0.542 \\
conv-AE      & agglom.   & 0.007 & 0.547 & 0.545 & 0.714 \\
conv-AE      & spectral  & 0.001 & 0.506 & 0.519 & 0.577 \\
conv-AE      & GMM       & 0.004 & 0.503 & 0.537 & 0.524 \\
conv-AE      & HDBSCAN   & 0.000 & 0.707 & 0.000 & n/a \\
hand-crafted & $k$-means & 0.557 & 0.823 & 0.901 & 0.554 \\
hand-crafted & agglom.   & 0.612 & 0.856 & 0.922 & 0.528 \\
hand-crafted & spectral  & 0.605 & 0.837 & 0.910 & 0.568 \\
hand-crafted & GMM       & 0.493 & 0.777 & 0.869 & 0.599 \\
hand-crafted & HDBSCAN   & 0.186 & 0.621 & 0.168 & 0.140 \\
HuBERT-ECG   & $k$-means & 0.133 & 0.591 & 0.708 & 0.572 \\
HuBERT-ECG   & agglom.   & 0.123 & 0.597 & 0.683 & 0.690 \\
HuBERT-ECG   & spectral  & 0.031 & 0.627 & 0.564 & 0.868 \\
HuBERT-ECG   & GMM       & 0.146 & 0.597 & 0.721 & 0.502 \\
HuBERT-ECG   & HDBSCAN   & 0.000 & 0.707 & 0.000 & n/a \\
ECGFounder   & $k$-means & 0.825 & 0.943 & 0.970 & 0.529 \\
ECGFounder   & agglom.   & 0.660 & 0.856 & 0.921 & 0.579 \\
ECGFounder   & spectral  & 0.801 & 0.930 & 0.964 & 0.537 \\
ECGFounder   & GMM       & 0.798 & 0.926 & 0.961 & 0.542 \\
ECGFounder   & HDBSCAN   & 0.261 & 0.626 & 0.226 & 0.202 \\
\midrule
\multicolumn{6}{l}{MIT-BIH ($n = 3\,000$, majority 0.500)} \\
beat-morphology & $k$-means & 0.652 & 0.875 & 0.933 & 0.517 \\
beat-morphology & agglom.   & 0.684 & 0.891 & 0.942 & 0.512 \\
beat-morphology & spectral  & 0.024 & 0.698 & 0.513 & 0.987 \\
beat-morphology & GMM       & 0.628 & 0.841 & 0.912 & 0.581 \\
beat-morphology & HDBSCAN   & 0.255 & 0.512 & 0.352 & 0.135 \\
conv-AE         & $k$-means & 0.121 & 0.598 & 0.666 & 0.708 \\
conv-AE         & agglom.   & 0.108 & 0.619 & 0.634 & 0.812 \\
conv-AE         & spectral  & 0.038 & 0.693 & 0.521 & 0.979 \\
conv-AE         & GMM       & 0.114 & 0.579 & 0.656 & 0.617 \\
conv-AE         & HDBSCAN   & 0.216 & 0.526 & 0.289 & 0.067 \\
raw-beat        & $k$-means & 0.018 & 0.560 & 0.567 & 0.734 \\
raw-beat        & agglom.   & 0.043 & 0.551 & 0.611 & 0.670 \\
raw-beat        & spectral  & 0.012 & 0.700 & 0.508 & 0.989 \\
raw-beat        & GMM       & 0.018 & 0.559 & 0.568 & 0.732 \\
raw-beat        & HDBSCAN   & 0.176 & 0.620 & 0.521 & 0.791 \\
\bottomrule
\end{tabular}
\end{table}

\bibliographystyle{elsarticle-num}
\bibliography{refs}
\end{document}
